% Page 095

\seite{95}

\abschnitt{1930}
\pstart
Ein neues Jahr hat begonnen. Es soll für uns\ob
ein Jahr der Freiheit werden. Die Rheinlande sollen\ob
endlich wieder frei werden von der feindlichen Be=\ob
satzung. Hoffen wir, daß dieser Wunsch in Erfüllung\ob
gehen möge! Möge das neue Jahr nur ein Jahr\ob
des unseren Freiheit werden — ein Jahr der Lösung\ob
und Ruhe!

\margmark{3. Januar    M}%
it Gottes Namen haben wir heute den Unter=\ob
richt wieder begonnen.\ob
Gestern fand zum zweitenmal die neue Gemeinde=\ob
und Wahl Aufstellung der Wahlausschüsse statt. Den\ob
Wahlvorstand sind gewählt: Da Walf, Heinz Hermann=\ob
dy\unsicher, Austen, Heinz Walf, Kath. Gerhard, Reinhard Peter\ob
und da Haupt Johann Krebs gegen die Gültigkeit der\ob
Wahl Einspruch erhoben war, wurde die Wahl\ob
mit 4 gegen 2 Stimmen für gültig erklärt. Die\ob
die sogenannten Dorf Gespannt\unsicher halten je 3\ob
Stimmen auf Heinands-Schütten und Schmitz Joh.\ob
Das lag da die Aufstellung bringen mußte,\ob
fiel auf Schmitz Gerhard der somit zum Gemeinde=\ob
vorsteher gewählt ist. Der Einspruch gegen die Wahl\ob
wird an den Kreisausschuß weitergereicht wer=\ob
den.

\margmark{11. Januar   E}%
inen wohlgelungenen Abend veranstaltete\ob
der Kriegerverein, Kameradschaft der Pfarrei\ob
Bettingen am Sonntag, den 5.\ Januar, für die Ange=\ob
hörigen seiner Mitglieder. Der mit den Landes=\ob
fahnen und den Bildern der Ehrungspräsidenten des
\pend
